\chapter{Introduction \& Background}\label{chap:introduction}

This chapter introduces the problem this dissertation addresses and the context in which it arises. It begins with the role of nanoscale structure and molecular state in organic materials, and with the reason conventional scattering and spectroscopy cannot measure both at a buried interface. It then traces how resonant soft X-ray methods developed to close that gap, from NEXAFS spectroscopy and microscopy through transmission scattering and reflectivity, and identifies the optical model used to interpret these measurements as the factor that now limits what they reveal. The chapter closes with the scientific questions and hypotheses that organize \crefrange{chap:znpc}{chap:xray-atlas} and an overview of the work.

\section{History and Background}\label{sec:history-and-background}

Organic materials describe a broad class of systems comprised mostly of carbon atoms bonded into complex molecular and polymer assemblies, including small molecules, polymers, peptides and proteins. Because of the rich bonding chemistry of carbon, this presents a zoo of molecular platforms that can be engineered and optimized around emergent phenomena set by both the innate bonding configurations and by the nano assemblies formed within the material due to processing conditions. For example in organic microelectronics such as solar cells and transistors, optimized bonding chemistries combined with thermodynamic insights on emergent nanostructures have produced devices bearing similar performance to traditional platforms but with novel optical, electric and kinetic properties in addition to being easier to manufacture at scale.\autocite{alqahtani2023green,collins2013absolute,alqahtani2018mixed,chen2024efficiency} Similarly, achievements have been made for lithography and patterning science leading to the development of novel photoresist chemistries and new emergent paradigms such as directed self-assembly pushing the final critical dimension to the nanometer limit.\autocite{adams2025polymer,sunday2026depth,prabhu2011photoresist} Additionally, polymer and peptide membranes are increasingly of interest to the development of new novel medical therapeutics focusing on both targeted drug delivery within nanoparticle polymeric carriers, and to unravel the complex lipid bilayers surrounding dangerous bacteria.\autocite{grabner2024molecular,ozturk2025highly}

The rich application of organic materials requires detailed knowledge of both the nano to mesoscale structures formed by the organic molecules and the molecular states such bonding configuration present within these structures. X-ray and neutron scattering measurements can uncover these hidden structures while providing some insights into the compositional gradients formed within them. Optimizing the size, shape, and composition within these structures has been critical to optimizing processing conditions around emergent function. However, with traditional single energy scattering it isnot generally possible to resolve the precise molecular states within these assemblies on their own. Because of this, molecular states and properties are usually determined separately through ultraviolet to visible spectroscopy, or X-ray microscopy. Yet the molecular state within nano structures, and especially at interfaces is often the most critical property that directly leads to the final emergent state. Tumbleston et al. showed that organic solar cell performance improves when the molecules responsible for charge generation orient face-to-face at donor acceptor interfaces\autocite{tumbleston2014influence}, with the same dependence of charge transport on molecular orientation since found in organic transistors\autocite{dindault2022correlating}. Neither the orientation nor the electronic state of molecules at such buried interfaces is accessible to the techniques above. \Cref{fig:1.1} places these scales in sequence, together with the scales that structural and spectroscopic measurements each reach.

\begin{figure}
  \centering
  \includegraphics[width=\linewidth,height=0.75\textheight,keepaspectratio]{hero/hero_fig1-01_scales}
  \caption[From technologies to the molecular state of a buried interface.]{From technologies to the molecular state of a buried interface. Organic materials enter devices, lithographic patterns, membranes, and thin films at the millimeter to micrometer scale. Their function depends on nanostructure near \qty{100}{nm}, on the interfaces buried within that nanostructure below \qty{10}{nm}, and on the orientation and electronic state of the molecules that form those interfaces below \qty{1}{nm}. The bar beneath marks the reach of conventional measurements. X-ray and neutron scattering and microscopy resolve structure, spectroscopy resolves electronic state, and a buried interface requires both at once (red dashed box).}
  \label{fig:1.1}
\end{figure}


\section{Resonant Scattering}\label{sec:resonant-scattering}

Resonant soft X-ray scattering (RSoXS) uses incident photon energy as a mechanism for contrast variation so that by changing energy the otherwise hidden molecular state gradients can be revealed.\autocite{grabner2025resonant} At energies near the ionization energy of core atomic electrons (1s, 2s, 2p, etc.) the bonding configurations between molecules produce a rich series of transitions into unoccupied molecular orbitals leading to the near edge X-ray absorption fine structure (NEXAFS).\autocite{stohr2010nexafs} \Cref{fig:2.4} shows an example of the NEXAFS differences between three polymer components. By tuning photon energies around these resonant differences, RSoXS reveals scattering differences dictated by the underlying quantum transitions.

This mechanism provides RSoXS with sensitivity to the chemical gradients within nanostructures that are otherwise missed by traditional scattering. Because of this, RSoXS characterization has been instrumental to the organic device and polymer physics community to unravel hidden nanostructures formed by the mixing properties of polymers that are otherwise unresolved by neutron and X-ray techniques. Critically, this technique is not limited to simple transmission scattering approaches but has seen wide utility across geometries including reflectivity and diffraction revealing chemical composition in otherwise inaccessible interfaces.

Chemical composition is only one mode of the chemical state, the molecular orientation and packing statistics are also critically relevant to emergent function, most notably dictating charge transport between molecules in electronic device applications. In conjugated materials the overlap of π orbitals between neighboring molecules sets the electronic coupling along each direction, so the same material transports charge efficiently in the film plane when its molecules stand edge-on and between stacked electrodes when they lie face-on.\autocite{sirringhaus1999two} In light-emitting diodes the orientation of the emitting transition dipoles likewise sets the fraction of generated light that escapes the device.\autocite{yokoyama2011molecular} This anisotropy leaves a fingerprint in the polarization dependence of the optical response itself, since the transition dipole moments responsible for each resonant feature couple most strongly to light polarized along their own orientation, the central idea developed in \cref{chap:molecular-films}.

\section{A Short History of Resonant Soft X-ray Characterization}\label{sec:a-short-history-of-resonant-soft-x-ray-characterization}

Resonant contrast in the soft X-ray range was first exploited for orientation rather than structure. Near edge X-ray absorption fine structure spectroscopy matured during the 1980s as a probe of molecules adsorbed on surfaces, where the polarization dependence of each resonance places the corresponding bond relative to the substrate.\autocite{stohr2010nexafs} Because the carbon, nitrogen, and oxygen K edges all lie in the soft X-ray range, the same measurement distinguishes chemical moieties in organic materials that are nearly indistinguishable to hard X-rays, and standardized methods for carbon K-edge spectroscopy of polymers and organic films followed.\autocite{watts2006methods} Scanning transmission X-ray microscopy at third-generation synchrotron sources\autocite{bilderback2005review} brought this chemical and orientational sensitivity into real space with resolution well below \qty{100}{nm},\autocite{kilcoyne2003interferometer,ade2008nexafs} and undulators delivering variable linear polarization made the polarization direction an experimental variable in its own right.\autocite{young2002variable} Resonant scattering developed in parallel at the transition metal L edges, where the magnetic contrast of 2p resonances resolved chiral magnetic domains in ultrathin films\autocite{durr1999chiral} and the resonant magneto-optical constants needed to interpret such measurements were measured directly.\autocite{kortright2000resonant} These two lines of development, chemically specific spectroscopy of organic matter and resonant scattering as a structural probe, met at the carbon K edge.

Carbon K-edge resonant scattering followed in 2006, when the first transmission RSoXS measurements at the Advanced Light Source resolved polymer nanoparticles in a chemically similar matrix without staining or labeling.\autocite{araki2006resonant,mitchell2006molecular} A dedicated instrument at beamline 11.0.1.2 followed in 2010 and, with real-time data processing, opened the technique to a broad user community. Over the following years RSoXS became the standard measure of domain size and domain purity in organic photovoltaic blends, and established quantitative correlations between scattering-derived morphology and device efficiency.\autocite{yan2012correlating,jiao2017quantitative} In 2012, Collins et al. showed that polarized RSoXS is sensitive to the local orientation of molecules through their transition dipole moments, revealing non-crystalline orientational order in organic solar cell blends.\autocite{collins2012polarized} Tumbleston et al. then correlated this interfacial orientation with device performance.\autocite{tumbleston2014influence} Polarization-dependent scattering from organic films has since been traced to the optical anisotropy of the constituent molecules.\autocite{gann2016origins}

The technique has since spread well beyond organic photovoltaics and beyond the carbon edge. A second dedicated facility for polarized RSoXS operates at NSLS-II\autocite{gann2021nist}, and measurements at the nitrogen K edge and at the tender sulfur K edge reach heteroatoms that carbon contrast cannot isolate\autocite{zhong2022decoupling,freychet2021resonant,mcneill2023resonant}. Applications now include the helical and cubic phases of liquid crystals\autocite{salamonczyk2017structure,cao2020deciphering,cao2020molecular}, proteins and plant cell walls\autocite{ye2018resonant,ye2018resonanta,rongpipi2021resonant}, the latent images of lithographic resists\autocite{freychet2019reconstructing,zhang2024characterization}, and in situ and operando studies of solutions and electrochemical interfaces\autocite{zhong2021probing,yang2022textless,grabner2024molecular}. GPU-accelerated simulation of polarized scattering from three-dimensional morphologies now allows such measurements to be modeled directly\autocite{saurabh2023textless}. Across this range, the interpretation of every measurement rests on the optical constants assumed for each component, the dependence examined in the remainder of this chapter.

\section{Resonant Soft X-ray Reflectivity}\label{sec:resonant-soft-x-ray-reflectivity}

Specular reflectivity measures the depth profile of the refractive index rather than its lateral structure. Parratt's recursion for the reflected amplitude of a stratified medium remains the basis of its analysis\autocite{parratt1954surface}, with interfacial roughness entering through the Névot-Croce factor\autocite{nevot1980caracterisation} and anisotropic layers treated exactly by 4 × 4 transfer matrix formalisms.\autocite{berreman1972optics,yeh1980optics,passler2017generalized} At hard X-ray energies, however, the indices of low-Z organic layers differ too little for their interfaces to be resolved. Wang, Araki and Ade showed that tuning to the carbon K edge restores this contrast\autocite{wang2005soft}, and resonant soft X-ray reflectivity was soon applied to buried interfaces in organic light emitting diodes and other organic devices\autocite{wang2008multilayer,hexemer2007d} and to molecular orientation in soft matter films.\autocite{mezger2011molecular} Dichroic resonant reflectivity of ultrathin anisotropic organic layers and of an organic single crystal then showed that the anisotropic optical constants of an organic film can be determined quantitatively from reflectivity alone,\autocite{capelli2019quantitative,pasquali2014structural,capelli2016quantitative} and resonant reflectometry has recovered site-selective spectra with depth resolution\autocite{hamannborrero2017site}.

Polarization adds orientation to this depth sensitivity. The optical tensor of an oriented molecular film is anisotropic, so s- and p-polarized reflectivity weight its in-plane and out-of-plane components differently, so that the depth profile of orientation can be resolved together with that of composition. Thelen et al. used polarized resonant reflectivity to measure orientation depth profiles in vapor-deposited organic glasses, and Ferron et al. resolved the orientation and conformation of molecules at the buried heterointerfaces of vapor-deposited films.\autocite{ferron2023composition,thelen2020molecular,ferron2022characterization} Each of these studies treats the optical constants as inputs, either tabulated from NEXAFS of a reference film or fit freely at each energy, with Kramers-Kronig consistency imposed as a constraint only when the energy sampling is dense enough to support it.\autocite{kuzmenko2005kramerskronig,stone2012kramers} The depth profile is then only as reliable as these optical constants, which is the limitation this work addresses.

\section{Optical Models and Their Limits}\label{sec:optical-models-and-their-limits}

Interpreting RSoXS requires the optical tensor of the material at every measured energy, the complex index of refraction \(n = 1 - \delta + i\beta\), whose dispersive part \(\delta\) sets the phase of the transmitted and scattered waves and whose absorptive part \(\beta\) sets their attenuation. Most approaches measure \(\beta(E)\) across the absorption edge by NEXAFS, recover \(\delta(E)\) through the Kramers-Kronig relations\autocite{yan2013accurate,watts2014calculation}, and assign each resonance to a molecular orbital and its transition dipole moment so that angle-resolved NEXAFS fixes the orientation dependence of the tensor.\autocite{stohr2010nexafs} This works when resonances trace cleanly to bonds of known dipole orientation, as for the π* transitions of planar aromatic rings. For the polymer N2200, however, Gann et al. showed that angle-resolved NEXAFS cannot trace the absorption features of a large, conjugated repeat unit to individual moieties.\autocite{gann2014near} Density functional theory supplies the missing assignment in principle, yet calculations routinely return thousands of transitions, place them on an absolute energy scale with an offset, broaden them by an imposed rather than computed width, and treat an isolated molecule, so they have reproduced measured spectra only qualitatively.\autocite{su2017first} Murcia developed a clustering approach that reduces a DFT calculation to a handful of principal transitions and refines them against measured spectra, reproducing the NEXAFS of phthalocyanines quantitatively while retaining the transition dipole of each cluster.\autocite{murcia2026quantitative} This dissertation builds directly on that framework.

Most work in RSoXS rests solely on this idea of contrast variation, choosing energies and polarizations that maximize contrast between pre tabulated species offering only qualitative insights into the chemical states in buried structures and at heterointerfaces. This presupposes that the potential molecules within these structures come from a set of pre characterized molecules that have been measured independently ahead of time. The natural consequence of this comes a lack of specificity against hydrocarbon contaminants and oxidations that may share similar spectral characteristics to the hypothesized compound family but deviate strongly outside of the few measurement energies. It is therefore never possible to exclude or characterize these contaminants. Similarly, at interfaces if novel molecular conformations or hybridization impact the optical response uniquely, then these characteristics cannot be uncovered. What is needed to resolve this is the full spectral measurement, collecting RSoXS across an entire absorption edge and with multiple polarizations. However, successfully analyzing this data requires a framework to produce accurate and facile anisotropic optical tensors that describe the energy and polarization dependence of the optical response, and are directly interpretable to uncover the molecular states responsible for this emergent response.

\section{Scientific Questions and Hypotheses}\label{sec:scientific-questions-and-hypotheses}

The framework called for above must supply an anisotropic optical tensor for every layer of a film at every energy and polarization. The reflectivity studies reviewed earlier obtain that tensor in one of two ways, tabulated beforehand from NEXAFS of a reference film and held fixed, or fit freely at each measured energy, with Kramers-Kronig consistency imposed only when the energy sampling is dense enough to support it. The first practice rests on the presumption that the molecules in every layer are the reference in some orientation and density, so a layer that departs from the reference can be neither interpreted nor explained. The second rests on no presumption about the molecule, so nothing about the molecule follows from it without a secondary analysis, and its parameter count grows with every energy added. One question is asked of each practice. The first concerns what becomes measurable when the tabulated tensor is replaced by one derived from the electronic structure of the molecule. The second concerns how the physical constraints that make such a derivation possible can be carried into the recovery of a tensor that no reference supplies, with the answer to the second resting on what the first establishes.

Murcia's clustering reduces a DFT calculation, which returns far more transitions than any spectrum resolves, to the few principal transitions that reproduce a measured NEXAFS spectrum once their strengths and widths are refined against it, with each retained cluster keeping the transition dipole assigned to it by the calculation. The tensor of a film is the ensemble average of that single-molecule response over the orientations and packing of the molecules within it, so the projection of each dipole through the orientation distribution of a layer gives the full energy and polarization dependence of that layer from its molecular tilt and density alone. The construction is not yet an analysis. The refinement is carried out against angle-resolved NEXAFS of a reference film whose own orientation enters the measured dichroism, so the dipole direction of each cluster and the tilt of the molecule are entangled in the same data used to fix both, and the clustering retains only the transitions selected as principal, so a weak, misassigned, or absent transition enters the tensor as an error present at every depth. Whether a tensor built this way can support a polarized resonant soft X-ray scattering analysis, and what such an analysis returns that the tabulated tensor cannot, is the first question of this work. \emph{\textbf{How can an optical tensor derived from ML-refined DFT be built into the analysis of resonant soft X-ray scattering, and what does that integration make measurable about the molecular state of a buried interface?}}

Much of what the integration would give can be reasoned out before it is attempted. With the spectral response fixed by the molecule, a layer is specified by its order and density, so the number of parameters no longer depends on how many energies are measured, and each added energy tightens the constraint on the same few parameters where it previously added its own set of tensor components to be fit. Because those parameters set the amplitude and phase of the reflectivity fringes at every energy and polarization, the energies that free fits treat as independent experiments become one overdetermined measurement of structure, with the thickness and sharpness of an interfacial layer constrained by how the fringes change across the edge in addition to their spacing at any one energy, so the features such a fit resolves need not be larger than the diffraction limit set by the largest momentum transfer measured. The dichroism of every resonance follows from the dipoles and the tilt, so the polarization dependence that a free fit absorbs into separate in-plane and out-of-plane components becomes a prediction. Because the tensor admits only the molecule, a fit can fail. A free tensor accommodates any layer, so a surface or buried layer whose optical constants differ from the bulk can be assigned to contamination or oxidation without that assignment ever being tested, whereas a tensor fixed to the molecule either reproduces such a layer with a different tilt and density or does not reproduce it, with the residual at the molecule's own resonances distinguishing the two.

What cannot be reasoned out in advance is how a wrong tensor behaves in a fit. A cluster carrying the wrong dipole, or a transition that the reference film did not contain, might appear as misfit at particular energies, or it might be absorbed into whichever structural parameters are free, biasing the recovered profile while the fit appears sound. The two cases can be separated only by fitting the same data with both tensors, a comparison that serves as the test of the first hypothesis. We hypothesize that existing reflectivity measurements contain more about the molecular state at buried interfaces than free-tensor analyses return, with a tensor derived from the molecule recovering that information and resolving order and density below the diffraction limit from fewer parameters and with smaller uncertainties than the free tensor requires. If the hypothesis holds, with a tensor that admits only the molecule reproducing the layers that free fits attribute to contamination, those layers become the molecule in another state, and the profile of orientation and density returned by the fit, together with the account of film growth read from it, becomes a product of this question, discussed where it arises without directing the work. If the hypothesis fails, the failure is itself informative, since the misfit then falls at the energies where the reference film and the buried molecules differ.

The premise of the first question is that the molecules in the film are the molecule of the reference. At an interface that premise can fail, since packing, conformation, and hybridization with a neighbor alter the unoccupied orbitals into which the core electron is excited. For species that cannot be prepared as standards, such as volatile solvents or molecules in conformations that occur only at an interface, it fails outright. The tensor must then be recovered from the scattering itself. Fitting it energy by energy restores the parameter explosion and discards a constraint that the physics supplies, since nothing in such a fit requires the optical constants found at one energy to be consistent with those found at another. The Kramers-Kronig relations tie the dispersion at every energy to the absorption at every other, so the dispersion below the edge, which sets the critical angles that anchor every reflectivity profile, is determined by the absorption across the whole edge. Nothing prevents optical constants fit independently from violating these relations, so such a fit can return dispersion that no causal medium could produce, with the violation undetectable short of sampling the entire edge. Any recovery of optical constants from scattering must therefore impose causality. Kuzmenko's variational analysis does so by expanding the optical response in Kramers-Kronig consistent basis functions, an approach that the triangle basis of Stone et al. has carried to resonant reflectivity across a full absorption edge, but each measured energy gains a basis function, so the dense sampling the method requires returns the parameter growth it was meant to remove. The alternative follows from the first question. The ML-DFT clusters are already a set of transitions at known energies with known dipoles. Written as oscillators whose line shapes carry their own Kramers-Kronig conjugates, they form a basis that is causal by construction, anisotropic through its dipoles, and fixed in size by the number of transitions in the molecule rather than the number of energies in the measurement. \emph{\textbf{How can the optical constants of a buried molecular state be recovered from the scattering itself, under the same physical constraints that made the first analysis possible, and what do they reveal about the electronic states that molecules adopt at an interface?}}

With such a basis the gains of the first question carry over. The parameter count is set by the number of resonances, so the strength of each resonance can be refined as a function of depth, and the orientation of its dipole separately from the molecular tilt, without the fit becoming underdetermined. The anisotropy of a buried layer then becomes a measured quantity, and because every trial spectrum is causal, no fit can return optical constants that a real medium could not possess. Reflectivity then returns a spectrum localized to a region that no reference measurement can isolate, a depth-resolved spectroscopy whose parameters trace back to families of electronic transitions. The interfacial states found under the first question, together with the observation that NEXAFS of films prepared under nominally identical conditions differs, leave open whether molecules within a monolayer of an interface occupy electronic states that differ from the bulk, with the oscillator basis allowing that question to be put to the data resonance by resonance. The measurement itself raises a second question, since exact optical constants at a few energies leave the spectrum between them unconstrained, so the choice between a few precisely measured energies and a continuous scan across the edge, settled so far by convention and by the cost of beamtime, acquires a quantitative answer. We hypothesize that an oscillator basis tied to transition dipoles allows the optical response of a buried region to be refined together with the structure, with the molecules at an interface differing from the bulk in the strengths of particular transitions, in ways that no rescaling of a reference spectrum reproduces, and with the spectrum recovered from a sparse measurement limited by how the edge is sampled before it is limited by the precision at any one energy. These claims are tested in Alq\textsubscript{3} and ZnPc against fits that hold a measured NEXAFS spectrum fixed, and on synthetic spectra sampled by several strategies. The limit in which no calculation supplies the transitions, and only their number is assumed, is discussed but remains to be tested on a measured interface.

Both questions rest on measured NEXAFS, from which the clusters are refined and against which recovered spectra are judged, yet spectra of nominally identical films differ. A spectrum must therefore be judged before it is reused, which requires that the raw measurement stay linked to every optical constant derived from it, that spectra from different beamlines share one normalization and processing convention, that quality be scored by reproducible metrics, and that every result carry its provenance under a versioned and citable identifier that analysis pipelines can reach programmatically. X-ray Atlas is built to these requirements, as the infrastructure both questions assume, with no question of its own.

\section{Overview}\label{sec:overview}

This work presents the development of facile anisotropic optical tensor models that directly connect the emergent response to both the physical molecular state and the electronic states within the molecules. Each resonant feature is traced to the family of electronic transitions that produces it. This builds on the developments of Dr. Victor Murcia's Ph.D. dissertation, which used machine learning to integrate density functional theory calculations of electronic states with experimental measurements of absorption spectroscopy.\autocite{murcia2026quantitative} Through this framework each resonant feature can be described by a handful of transitions, each bearing a transition dipole moment that describes its anisotropic response. This isolates the single molecule optical response so that each molecular assembly can be constructed from the correct ensemble averaging of the molecular populations within it, allowing the direct measurement of properties such as molecular orientation and packing without a secondary analysis. Applying this to RSoXS in reflection geometry constrains the molecular state evolution far below the nominal diffraction limit.

Because each cluster carries its own oscillator strength and transition dipole, the resonances themselves become quantities that can be tracked through a film. Where the molecular state at a heterointerface alters the electronic structure, through a new conformation, hybridization with a neighbor, or a change in packing, the amplitudes and dipole orientations of the affected resonances change with depth while the others do not. Refining them as functions of depth recovers a spectrum localized to regions that no reference measurement can isolate, such as a buried interface a single molecular layer thick. This approach, developed in \cref{chap:resonances}, extends resonant reflectivity from contrast variation to depth-resolved spectroscopy.

Finally, this dissertation presents a modern platform for collaborative open data surrounding anisotropic optical tensors. Programmatic, standardized access to reusable optical models for molecular materials is critical to tracing the molecular evolutions within a film. Present best practices make data `available upon request', yet this denies the programmatic access that cross-publication meta-analyses require, leaving unexamined the phenomena that each analysis misses on its own. This part of the work is a collaboration among seven institutions in three countries.

The remainder of this dissertation is organized around these questions. \Cref{chap:molecular-films} sets out the materials and optical background of the work. It describes the device properties that molecular packing controls, the growth of vapor-deposited ZnPc and Alq\textsubscript{3} films and the open problems of their first molecular layer, the core excitations that give each resonance a transition dipole and the contrast it produces, the oscillator and Kramers-Kronig relations that turn those resonances into optical constants, and the DFT calculations, clustering, and refinement that supply the resonances of each molecule. It closes by projecting the refined molecular tensor into the frame of a film and testing a film of uniform tilt against angle-resolved NEXAFS of ZnPc. \Cref{chap:methods} contains the fundamentals and methods of polarized resonant soft X-ray reflectivity used throughout. Far-field scattering is recast around the specular geometry, with the reason that specular reflectivity cannot be treated kinematically made explicit, and the orientation-gradient profile used in \cref{chap:znpc} is derived. The measurement at ALS beamline 11.0.1.2 is described from sample alignment through scan configuration, reduction, stitching, and the uncertainty budget. The analysis methods cover the exact transfer matrix calculation for stratified anisotropic media, the interfacial roughness models, objective transforms, optimizers, posterior estimation, and model selection, with Bayesian optimization and amortized inference presented as the route to autonomous measurement.

In \cref{chap:znpc}, which reproduces the ZnPc manuscript, the first question is put to vapor-deposited ZnPc films on native-oxide silicon. The single-molecule tensor is constructed from DFT transition clusters refined against angle-resolved NEXAFS of an edge-on film, and the film is modeled first as three slabs, each described by a density and a molecular tilt, against the conventional alternative in which the uniaxial tensor components are fit freely at every energy in every slab. The ML-DFT tensor reproduces the reflectivity across the pre-edge and the resonance series with equal or smaller residuals from an order of magnitude fewer parameters, with correspondingly smaller uncertainties. Because the tensor admits only ZnPc, the contamination or oxidation customarily assigned to interfacial layers is not required by the data within the sensitivity of the resonance series. The slab structure is then replaced by a single continuous profile in which orientation and density relax from a pinned state at each interface into the bulk over a characteristic length. This profile fits better with fewer parameters still, and with it a face-on, lower-density state is resolved within one molecular layer of the oxide, along with a denser state oriented near the magic angle at the free surface, both finer than the diffraction limit. These profiles are read against substrate patterning and against a growth front on which substrate adhesion competes with π-π cohesion, with the consequence that coverage-checkpoint studies record a transient state of the growth front rather than the buried interface of a continuously deposited film.

In \cref{chap:resonances} the second question is put to two systems. The optical response is written as a sum of Kramers-Kronig consistent oscillator pairs, each a Gaussian absorption with its Dawson-function conjugate, tied to the ML-DFT transition clusters so that every oscillator carries a transition energy, a width, an oscillator strength, and a dipole orientation from which the tensor follows. For Alq\textsubscript{3}, whose NEXAFS differs between films prepared under nominally identical conditions, the amplitudes of the principal resonances are refined with depth, removing a degeneracy between the optical model and the film density that a borrowed spectrum can resolve only by driving the substrate density to an unphysical value. For the ZnPc films of \cref{chap:znpc}, the amplitude and dipole cone angle of each principal π* resonance are refined and compared with the conventional NEXAFS-fixed tensor, with the improvement concentrated at the energies of strongest orientation sensitivity and with the two lowest resonances exchanging oscillator strength between the free surface and the oxide interface. A synthetic carbon edge of known transitions is then used to compare sampling strategies, energies placed on the resonances, drawn at random, or chosen by a Gaussian-process design, against a continuous scan with larger noise at each point. Sparse measurements are found to be limited by the number and placement of energies before they are limited by the precision at any one of them. The same construction serves as a generator of Kramers-Kronig consistent optical constants for planning measurements.

\Cref{chap:xray-atlas} presents X-ray Atlas, an open and programmatic database of NEXAFS spectra and of the optical constants and versioned optical models derived from them, built around the provenance, quality metrics, and citable identifiers that reuse requires. The data model links each raw measurement to its processed spectrum and derived optical constants, every contribution is normalized under one server-side convention and scored by fixed metrics of energy resolution, signal to noise, and normalization residual, contributors are identified through ORCID, and each processing revision receives a persistent DOI. A REST interface, a Python package, and a command-line client extend the same records to scripted analysis. Optical models pin a dataset revision, so a tensor such as the one of \cref{chap:znpc} can be forked, cited, and reproduced. The answers to both questions are collected in \cref{chap:conclusion}, together with the extensions they open, finite element and neural surrogate solvers for geometries without an exact layered solution, the oscillator description applied to polymers and other absorption edges, and autonomous measurement at the beamline.
